% GENERATED from content/data/enrichment.json by mxbuild.enrichment; do not edit.
\mxunit{enrichment-lattice-paths-notes}{Chemins, combinaisons et Pascal}
\mxsubtitle{Première--Terminale · dénombrement \textperiodcentered{} Leçon et méthode \textperiodcentered{} 05}

\begin{prerequisites}
Produit, factorielle et choix sans ordre.
\end{prerequisites}

\begin{keypoint}[Idées et méthode]
Un chemin de $(0,0)$ à (a,b) n’utilisant que des pas vers la droite ou vers le haut comporte $a + b$ pas. Choisir les positions des a pas horizontaux donne $C(a+b,a)$ chemins.

Pour $0 \leq  k \leq  n$, $C(n,k) = \frac{n!}{k!(n-k)!}$. L’ordre des objets choisis ne compte pas. La symétrie $C(n,k) = C(n,n-k)$ revient à choisir les objets retenus ou ceux exclus.

La relation de Pascal $C(n,k) = C(n-1,k-1) + C(n-1,k)$ partage les choix selon la présence d’un objet fixé. On peut aussi classer les chemins selon leur dernier pas.
\end{keypoint}

\begin{example}
De $(0,0)$ à $(3,2)$, il y a $C(5,3) = 10$ chemins. Pour passer par $(1,1)$, il faut choisir un des deux chemins jusqu’à ce point, puis un des trois chemins jusqu’à l’arrivée : $2 \times  3 = 6$ chemins.
\end{example}

\begin{figurebox}
\begin{tikzpicture}[scale=.85]\draw[mxline](0,0)grid(3,2);\draw[mxmint,very thick,->](0,0)--(1,0)--(1,1)--(2,1)--(3,1)--(3,2);\fill[mxink](1,1)circle(3pt)node[above left]{$(1,1)$};\node[below]at(0,0){$(0,0)$};\node[above]at(3,2){$(3,2)$};\end{tikzpicture}
\end{figurebox}

\begin{activity}[À explorer]
Dessiner un petit quadrillage et écrire dans chaque sommet le nombre de chemins qui y arrivent. Additionner les valeurs du voisin de gauche et du voisin du bas.
\end{activity}

